% Finalizado 0 
\begin{center}
    \textbf{Abstract}
\end{center}

\noindent This Final Degree Project presents a graphics engine implemented using Vulkan. The primary objective is to achieve high performance in real-time particle simulation and rendering. To this end, a parallel architecture based on \textit{Compute Shaders} has been designed, enabling the simulation of complex volumetric effects such as fire, smoke, and water. The cornerstone of the system is the complete delegation of the particle lifecycle to the GPU. Consequently, the main processor is relieved from updating positions, photometric properties, and physical dynamics. This approach effectively mitigates the traditional data transfer bottleneck between the CPU and video memory, a common limitation in legacy or high-level graphic interfaces.

The graphics engine is predicated on exhaustive hardware control, structured through independent modules. The pipeline integrates procedural computation, explicit synchronization via memory barriers (\texttt{VkBufferMemoryBarrier}), and volumetric rendering utilizing \textit{Point Sprites} in conjunction with \textit{Alpha Blending}. In addition to the particle systems, the application provides a three-dimensional environment equipped with an empirical ADS (Ambient, Diffuse, and Specular) lighting model, an optimized geometry parser, and a dynamic \textit{Skybox}. The utilization of descriptors and local buffers ensures low-latency resource access. Experimental results corroborate Vulkan's efficiency: the system is capable of concurrently processing a large quantity of particles, sustaining a stable rate of 60 frames per second on mid-range graphic processors while offering extensive visual parameterization options. In conclusion, the project constitutes a robust and scalable architectural foundation for integrating physical simulations into contemporary graphical applications.

\vspace{0.5cm}
\textit{Keywords:} vulkan, compute shaders, particle systems, real-time, gpu

\vspace{0.5cm}
\textit{Specific competencies:} 
In this work, the following specific competencies have been developed:
\begin{itemize}
    \item Ability to design, program, and evaluate interactive systems that display complex information. Practical application of this knowledge to solve real-world problems within human-computer interaction design.
\end{itemize}